% ===== §1 Introduction =====
\section{Introduction}
\label{sec:intro}

\noindent
Most of what is ever perceived as beautiful is perceived in the course of an ordinary day, and almost none of it is perceived in a museum. A person wakes to a particular quality of light on a wall, hears rain begin against the glass, tastes the first mouthful of something warm, watches someone they love cross a room. Such perceptions are small in amplitude, high in frequency, and gone almost before they are noticed, and philosophy has largely declined to treat them, having inherited a discipline of aesthetics built around the rare and framed object. A theory of beauty attending only to what hangs in galleries has thereby passed over the overwhelming majority of the sites at which the aesthetic actually occurs. This paper takes those sites for its subject. It concerns the aesthetics of the everyday, and it concerns what it would be to cultivate a capacity for that aesthetics, which is the old question of aesthetic education asked of the most ordinary material there is.

The everyday is not an easy subject, and it is not easy for the reason that recommends it. The same low amplitude and high frequency that make daily aesthetic life pervasive also make it the natural home of habituation: the light on the wall becomes \emph{the light on the wall}, unseen by the second week. An account of everyday aesthetics that only celebrated the beauty available in ordinary things, without reckoning with the standing tendency of ordinary things to go dead to perception, would be sentimental and would deserve the neglect the subject has received. The everyday is the hardest case for any claim that the aesthetic is generated rather than given, because it is where the givenness of the world is most nearly total and where the work of generating perception must be done against the heaviest resistance. This paper accepts that difficulty as its proving ground: a theory of aesthetic generation that can be made to hold in the kitchen at seven in the morning has been tested where it is most likely to fail.

\subsection*{The proposal: generative relational aesthetic education}

The central proposal can be put as one thesis with four faces, and the argument of the whole paper is the defence of these four together. The first is ontological, and it is the ground of the other three. Beauty is generated within a relational structure. It is not projected onto a neutral object by a sovereign subject, in the manner of aesthetic subjectivism, nor discovered as a property the object carried in itself waiting to be found, in the manner of aesthetic objectivism; it comes to be in the relation between a perceiver and what is perceived, and, in the cases this paper cares about most, in the relation between two perceivers turned together toward a shared world. A meal, a swept room, the light at a particular window: the beauty of these is not located in the thing, nor in either mind alone, but in the relational event of \emph{our} perceiving and \emph{our} doing thus, together. The everyday is where this is most visible, because everyday things are too plain to sustain the illusion that their beauty was simply sitting in them.

The second face follows from the first and gives the paper its title. If beauty is generated in a relational structure, then an education in beauty must cultivate the relation, and not only the individual within it. This marks the proposal off from the tradition of aesthetic education it most respects. From Schiller's letters onward \citep{schiller1967}, that tradition has taken its unit to be the single person: it asks how one human being's sensibility, taste, or character might be brought toward completion, how the sensuous and the rational might be reconciled within a soul. The proposal here asks a different question, whose difference is the whole contribution: how a relation might be brought to go on generating the aesthetic, how two people might cultivate between them a shared capacity that neither carries alone. The cultivated thing is the relation's power to generate perception, and this is why the paper calls its subject generative relational aesthetic education.

The third face guards the second against a misreading it invites. To make the relation the unit of cultivation is not to dissolve the individual into it. The relational and the individual stand in a dialectic, and the dialectic is a spiral rather than a subordination. An individual's perceptual refinement, the trained discrimination this paper will trace through the neuroscience of expertise, is the condition on which relational aesthetic generation draws; and the shared perception generated in the relation returns to reshape the individual's own capacity, which enters the next round already altered. Neither pole is prior. The individual sinks into the relation and the relation returns to remake the individual, and the return does not arrive back at its origin. This spiral, which earlier papers of this series have tracked under the geometry of holonomy, is what keeps the relational from swallowing the individual, and the paper will insist on it wherever the temptation to a simpler collectivism appears.

The fourth face concerns the bearing of all this on ethics, and it is placed here, at the outset, because it is not an addition to the aesthetic argument but a consequence of it. Because beauty is generated in a relation, and because a relation may be of a kind that returns what it generates to both who stand in it or of a kind that draws one of them off as mere material for the other's experience, aesthetic generation carries an ethical valence at its root. A perception that consumes the other as a spectacle, the aestheticising gaze that turns a person into scenery, is a bad aesthetic generation in the exact structural sense this series has given to badness; a perception in which each party to the shared seeing is a beneficiary of it is a good one. Beauty and ethics are not joined here by adding a moral rule to an aesthetic fact. They share a single root in the structure of relational generation, and the theory of justice this series has been building enters the aesthetics of the everyday through that shared root rather than over a wall.

\subsection*{Contribution, and its measure}

The contribution of this paper should be stated with its limits showing. It does not offer a completed system of relational aesthetic education, and it does not claim to have settled how such an education would proceed. What it offers is, first, the rethinking of aesthetic education along three axes on which that tradition has rarely been set: the relational system, in which the unit of cultivation is a relation rather than a person; the theory of justice, in which the uneven distribution of perceptual capacity is recognised as a form of injustice; and ethics, in which aesthetic generation is shown to carry its ethical valence at the root rather than by later addition. On the ground so prepared, the paper then proposes, as a preliminary and openly revisable offering, four models of everyday aesthetic practice, each with its theoretical descent made explicit and none advanced as a closed prescription. That the models are four, and offered for testing and amendment rather than as a finished inventory, is not a modesty of manner but a fidelity to the subject: an aesthetic education that could be reduced to a fixed procedure would have contradicted, in its form, the generativity it means to cultivate.

\subsection*{Where the paper stands, and how it proceeds}

The argument does not begin from nothing. It stands on ground this series has already laid: a relational ontology in which relations precede their terms; a theory of value in three registers, imaginary, symbolic, and real, whose real register alone carries endogenously good value; the distinction between the good and bad cycle, given a geometric body as the sign of a holonomy; the framework of generative justice developed in exchange with Eglash \citep{eglash2016}, in which every party to a generative process must be a beneficiary and none reduced to fuel; and the account of happiness, from an earlier paper of this series, as a quality of a dynamical process rather than a state to be reached. To these it joins resources from outside: the phenomenology of perception, the predictive-coding account of the perceiving brain, the psychoanalytic account of how desire organises the visual field, the political economy of the historically produced senses, the tradition of everyday aesthetics, and Ranci\`ere's account of the distribution of the sensible. The joining is deliberate and is governed by a method the paper states plainly, so that its breadth is not mistaken for eclecticism. The disciplines are arranged in three tiers. A small number bear the weight of the argument's spines. Several more enter as structural interventions, each doing a determinate piece of work at a determinate place. The rest sound as overtones, supplying an image or a formal correspondence without being asked to carry a claim. Where the paper reaches a formal register at all, it keeps the formalism in the serving position, translating results already established rather than deriving them here; and throughout, it declines to convert the cultivation of perception into a technique, on the methodological ground that a generativity made procedural is a generativity killed.

Three spines carry the paper, and each is named here with the disciplines that bear it. The first is the spine of \emph{generation}: that aesthetic perception generates a world rather than registering one, a claim tested where it is hardest, in the everyday, and argued across phenomenology, cognitive science, and psychoanalysis. The second is the spine of \emph{the body}: that perception is plastic, so that the cultivation of perception is a real training and not a metaphor, a claim given its empirical body by the neuroscience of perceptual learning. The third is the spine of \emph{justice}: that the distribution of perceptual capacity is uneven in ways that constitute a neglected injustice, so that aesthetic education is answerable to the theory of justice, a claim argued through the political economy of perception and the distribution of the sensible. The paper is the working-out of these three spines toward the proposal they were raised to support, and it ends not in a synthesis but in a set of directed fissures, each marking a place where what has been said runs out and points beyond itself, the first of them toward the emergence of culture that a later paper must take up.
